% SPDX-License-Identifier: GPL-3.0-or-later
% Copyright (C) 2026 Alessandro Greco
\chapter{Эксплуатация реле или моста}\label{cap:22-relay}
\aggiornatoal{23 сентября 2026 г.}{tor 0.4.9.13 (руководство); лабораторная сеть}

\begin{obiettivi}
К концу этой главы вы сможете:
\begin{itemize}[nosep]
  \item выбирать между промежуточным реле, выходным реле и мостом с учётом доступных ресурсов и приемлемых рисков;
  \item писать базовую конфигурацию для каждого варианта, включая ограничения пропускной способности и хранение ключей вне сети;
  \item читать и составлять политику выхода и проверять её влияние в лабораторной сети (Лабораторная работа~14);
  \item правильно объявлять семейство реле.
\end{itemize}
\end{obiettivi}

Сеть Tor существует благодаря тысячам людей и организаций, предоставляющих пропускную способность и компьютеры. В предыдущих главах мы рассматривали реле с точки зрения клиента; теперь перейдём на другую сторону. Эта глава не является руководством по системному администрированию — Tor Project поддерживает актуальное руководство оператора, — а объясняет, что означают основные эксплуатационные решения на уровне протокола.

% ---------------------------------------------------------------------------
\section{Какой тип узла выбрать?}\label{sec:relay-tipi}

\begin{description}
  \item[Промежуточное реле] (которое со временем может также стать входным реле, guard): пересылает ячейки между другими реле, но никогда не открывает соединения к внешним интернет-ресурсам. Это наименее требовательная роль. Его адрес присутствует в консенсусе, но веб-сайты не видят трафик, исходящий от него.
  \item[Выходное реле] Открывает соединения к назначениям от имени клиентов. Эта роль особенно ценна, поскольку выходных реле мало относительно спроса (\cref{sec:percorso-pesi}), но именно она создаёт наибольшую внешнюю нагрузку и риск: получатели вредоносного или нежелательного трафика видят IP-адрес выходного реле как кажущийся источник.
  \item[Мост] Это реле, не включённое в публичный список и используемое как первый переход людьми, которые не могут подключиться к публичной сети (\cref{sec:cens-bridge}). Даже скромное домашнее соединение может быть полезно в сочетании с подключаемым транспортом, например obfs4 или WebTunnel.
\end{description}

\begin{attenzione}[Перед началом]
Эксплуатация выходного реле имеет юридические и операционные последствия, включая сообщения о злоупотреблениях и запросы от органов власти; они рассматриваются в \cref{cap:25-legale}. Разумно начать с промежуточного реле или моста, уведомить провайдера связи и не размещать реле на том же интернет-подключении, которое используется для чувствительной деятельности через Tor в роли клиента.
\end{attenzione}

% ---------------------------------------------------------------------------
\section{Базовая конфигурация}\label{sec:relay-config}

Для промежуточного реле достаточно нескольких строк \file{torrc}:
\begin{lstlisting}[numbers=none]
Nickname        MyRelay
ContactInfo     operator@example.org
ORPort          443
ExitRelay       0
SocksPort       0
RelayBandwidthRate  4 MBytes
RelayBandwidthBurst 8 MBytes
\end{lstlisting}
\opt{ContactInfo} позволяет другим связаться с оператором при проблемах с реле. Руководство напоминает, что дескрипторы архивируются и публикуются, поэтому этот адрес может оставаться публично доступным неопределённо долго. \opt{ExitRelay 0} явно фиксирует поведение, которое уже является значением по умолчанию в современных версиях: при значении \texttt{auto} tor становится выходным реле только тогда, когда оператор задаёт один из параметров, разрешающих выходной трафик. \opt{RelayBandwidthRate} и \opt{RelayBandwidthBurst} ограничивают среднюю и пиковую пропускную способность. В отличие от них, \opt{AccountingMax} ограничивает общий объём трафика за расчётный период, например месяц. После достижения лимита реле перестаёт принимать новый трафик до следующего периода.

\textbf{Практическое замечание.} При значении \texttt{ExitRelay auto} к включению режима выходного реле могут привести не только явные правила \texttt{ExitPolicy}, но и настройки \texttt{ReducedExitPolicy} или \texttt{IPv6Exit}. Напротив, \texttt{ExitRelay 0} из примера полностью запрещает выходной трафик. Перед запуском публичного реле проверяйте эффективную конфигурацию.

\subsection{Жизненный цикл нового реле}

Только что запущенное реле публикует свой дескриптор авторитетным серверам каталога, которые добавляют его в консенсус и начинают измерять (\cref{sec:dir-voto}). Сначала оно получает мало трафика: клиенты выбирают реле пропорционально весам их пропускной способности, которые остаются низкими, пока измерения не стабилизируются. Флаги присваиваются со временем (\cref{sec:dir-flag}). Для \texttt{Stable} требуется длительная непрерывная работа, а для \texttt{Guard} дополнительно необходимы достаточный возраст и высокая доступность. После получения флага Guard объём трафика может сначала даже уменьшиться: клиенты, уже выбравшие свои входные реле, не начинают использовать новый узел как промежуточный с той же частотой. Затем трафик постепенно растёт по мере того, как новые клиенты выбирают его (\cref{sec:guard}).

Два вида сбоев не эквивалентны. Частые перезапуски без смены идентичности сохраняют измерения пропускной способности, но уменьшают непрерывную доступность, необходимую для \texttt{Stable} и \texttt{Guard}; эти флаги приходится заслуживать заново. Смена же \emph{идентификационных ключей} создаёт в представлении authorities совершенно новое реле, начинающий с нуля даже по измерениям пропускной способности.

\subsection{Ключи}

Идентичность реле определяется его основным ключом Ed25519 (\cref{sec:canale-chiavi}). По умолчанию tor хранит его на диске реле. При \opt{OfflineMasterKey 1} основной ключ остаётся на другом компьютере; на сервер передаются только временные ключи подписи, сертифицированные с помощью \cmd{tor --keygen}. Если сервер будет скомпрометирован, атакующий получит истекающие временные ключи, а не долгосрочную идентичность реле. Этот подход требует дисциплины: временные ключи нужно обновлять до истечения срока действия, иначе relay перестанет работать.

% ---------------------------------------------------------------------------
\section{Политика выхода}\label{sec:relay-exitpolicy}

Политика выхода — упорядоченный список правил \texttt{accept} и \texttt{reject}, каждое из которых задаёт адрес или сеть и порт либо диапазон портов. Tor обрабатывает правила \textbf{по порядку}, применяя первое совпавшее. Правила оператора располагаются перед правилами по умолчанию, перечисленными в руководстве. По умолчанию блокируются порты электронной почты (25), NNTP (119), NetBIOS и SMB (135--139, 445), а также некоторых приложений обмена файлами; остальной трафик разрешается. Отдельно tor по умолчанию отклоняет приватные назначения, включая \texttt{127.0.0.0/8}, \texttt{10.0.0.0/8} и \texttt{192.168.0.0/16}, а также собственный публичный адрес реле. Это защищает внутреннюю сеть оператора (\opt{ExitPolicyRejectPrivate}). В лабораторной сети эта защита отключена, чтобы разрешить трафик к \texttt{127.0.0.1}.

\textbf{Порядок защитных правил.} Если включён \texttt{ExitPolicyRejectPrivate} (по умолчанию), запреты на приватные сети и локальные адреса добавляются \emph{перед} правилами оператора. В частной сети Chutney эта защита намеренно отключена; такую настройку нельзя переносить на публичное выходное реле.

Выходное реле, ограниченное веб-трафиком, может использовать:
\begin{lstlisting}[numbers=none]
ExitRelay  1
ExitPolicy accept *:80
ExitPolicy accept *:443
ExitPolicy reject *:*
\end{lstlisting}
Tor Project также рекомендует \opt{ReducedExitPolicy} как распространённый компромисс: она разрешает набор широко используемых сервисных портов и исключает порты, чаще вызывающие жалобы на злоупотребления.

\begin{figure}[htbp]
\centering
% SPDX-License-Identifier: GPL-3.0-or-later
% Copyright (C) 2026 Alessandro Greco
%% Русское альтернативное описание EPUB.
\altfig{Политика выхода узла Chutney: только в тестовой сети ExitPolicyRejectPrivate 0 отключает предварительное отклонение приватных адресов. Два запроса BEGIN к 127.0.0.1 на порты 8080 и 8081 проходят правила в порядке записи. Порт 8080 совпадает с reject 127.0.0.1:8080, результат END EXITPOLICY. Порт 8081 пропускает это правило и совпадает с accept 127.0.0.0/8:*, что разрешает попытку TCP-подключения; CONNECTED возможен лишь при работающем сервере. Если поменять правила местами, запрет не сработает. На публичном выходном реле отключать ExitPolicyRejectPrivate нельзя.}
% Политика выхода только в лаборатории
\begin{tikzpicture}[font=\sffamily\scriptsize,node distance=0.45cm,
  reg/.style={draw,rounded corners=1pt,minimum height=0.6cm,text width=4.6cm,font=\ttfamily\scriptsize},
  esito/.style={draw,rounded corners=2pt,minimum height=0.6cm,text width=2.5cm,align=center}]
  \node[font=\sffamily\scriptsize] (top) at (0,1.25) {два BEGIN: 127.0.0.1:8080 / :8081};
  \node[draw,rounded corners=2pt,fill=torviola!10,font=\ttfamily\tiny] (r0) at (0,0.15) {ExitPolicyRejectPrivate 0};
  \node[reg,fill=avviso!12,below=of r0] (r1) {reject 127.0.0.1:8080};
  \node[reg,fill=torverde!12,below=of r1] (r2) {accept 127.0.0.0/8:*};
  \node[reg,fill=torgrigio!10,below=of r2] (r3) {(прочие правила по умолчанию)};
  \draw[->] (top) -- (r0); \draw[->] (r0) -- (r1); \draw[->] (r1) -- node[right,font=\sffamily\tiny] {8081: не совпадает} (r2); \draw[->] (r2) -- node[right,font=\sffamily\tiny] {иначе} (r3);
  \node[esito,fill=avviso!15,draw=avviso,right=1.2cm of r1] (e1) {\texttt{END} причина\\\texttt{EXITPOLICY}};
  \draw[->,avviso,thick] (r1) -- node[above,font=\sffamily\tiny] {8080 совпадает} (e1);
  \node[esito,fill=torverde!15,draw=torverde!60!black,right=1.2cm of r2] (e2) {попытка TCP\\CONNECTED лишь\\при сервере};
  \draw[->,torverde!60!black,dashed] (r2) -- node[above,font=\sffamily\tiny] {8081 совпадает} (e2);
  \node[anchor=west,font=\sffamily\tiny,text=torgrigio,align=left] at (-5.9,-0.55) {первое совпавшее\\правило решает};
  \node[anchor=west,font=\sffamily\tiny,text=avviso] at (-5.9,-1.95) {не для публичного exit};
  \node[font=\sffamily\tiny,text=torgrigio,anchor=south] at (r0.north) {только тест: ExitPolicyRejectPrivate 0};
\end{tikzpicture}

\caption{В лаборатории Chutney \opt{ExitPolicyRejectPrivate 0} включается только для тестов с приватными адресами; обычно Tor отвергает их до правил оператора. \texttt{BEGIN} к \texttt{127.0.0.1:8080} совпадает с \texttt{reject} и приводит к \texttt{END EXITPOLICY}; порт 8081 проходит это правило и совпадает с \texttt{accept}. Разрешение позволяет попытку TCP-соединения, но \texttt{CONNECTED} требует работающего сервера. Если поменять правила местами, запрет не подействует.}
\label{fig:policy}
\end{figure}

\paragraph{Что получает клиент.} Дескриптор реле содержит полную политику, а консенсус и микродескрипторы — только сводку по портам, например \texttt{p accept 80,443}. Клиент использует её, чтобы решить, способно ли выходное реле обслужить запрос (\cref{sec:percorso-vincoli}). Поскольку адреса в этой сводке отсутствуют, о запрете, зависящем от адреса, клиент иногда узнаёт только после попытки соединения, как показывает лабораторная работа.

% ---------------------------------------------------------------------------
\section{Лабораторная работа: проверка политики выхода}\label{sec:lab-relay}

\begin{laboratorio}[14]{Правило выхода и его последствия}
\textbf{Цель.} Добавить правило к лабораторному выходному реле и с помощью \file{labs/ru/код/torlab} наблюдать ответы на разрешённые и запрещённые назначения.\\
\textbf{Предварительные условия.} Лабораторные работы~3 и~4; лабораторная сеть должна быть запущена.\\
\textbf{Оценочное время.} 1 час.
\end{laboratorio}

\paragraph{Шаг 1: правило.} В \file{torrc} узла \texttt{test006e} вставьте правило, запрещающее один порт, перед существующими правилами, разрешающими \texttt{127.0.0.0/8}:
\begin{lstlisting}[numbers=none]
ExitPolicy reject 127.0.0.1:8080
ExitPolicy accept 127.0.0.0/8:*
\end{lstlisting}
Листинг показывает порядок двух соответствующих правил: в существующем \file{torrc} достаточно вставить новую строку \texttt{reject} перед первой строкой \texttt{ExitPolicy}, не дублируя имеющийся \texttt{accept}. Порядок имеет значение. Если поместить \texttt{reject} после \texttt{accept}, запрет не сработает, поскольку побеждает первое совпавшее правило. Именно такая ошибка была зафиксирована в предыдущем эталонном эксперименте. После перезапуска выходного реле запустите два небольших HTTP-сервера на портах контейнера 8080 и 8081.

\paragraph{Шаг 2: три запроса.} Используя \texttt{build} и \file{stream.py}, постройте цепь, заканчивающуюся на \texttt{006e}, и отправьте три сообщения \texttt{BEGIN} к портам 8080, 8081 и 9999, где ничего не слушает. Ответы в одном из наших сеансов (коды, адрес и TTL могут отличаться в зависимости от версии и среды):
\begin{lstlisting}[numbers=none]
127.0.0.1:8080 -> END  reason=EXITPOLICY address=127.0.0.1 TTL=1800
127.0.0.1:8081 -> CONNECTED
127.0.0.1:9999 -> END  reason=MISC
\end{lstlisting}

\paragraph{Шаг 3: интерпретация ответов.} Отказ по политике приводит к причине \texttt{EXITPOLICY}. Ячейка \texttt{END} также содержит разрешённый адрес и TTL. Клиент может запомнить, что данное выходное реле не обслуживает это назначение, и выбрать другой, не дожидаясь нового дескриптора. В отличие от этого, порт без слушающего сервера возвращает общий \texttt{MISC}, а не \texttt{CONNECTREFUSED}, который спецификация перечисляет для отказа в соединении~\spec{tor-spec}{Closing streams}.

\begin{approfondimento}[Почему \texttt{MISC}, а не \texttt{CONNECTREFUSED}?]
Ответ находится в \texttt{info}-журналах выходного реле. Прослеживание этой ситуации — хороший пример исследования расхождения между спецификацией и реализацией:
\begin{lstlisting}[numbers=none]
[info] connection_handle_write_impl(): in-progress connect failed.
       Removing. (Connection refused)
[info] errno_to_stream_end_reason(): Didn't recognize errno 115 (Operation
       now in progress); telling the client that we are ending a stream
       for 'misc' reason.
\end{lstlisting}
Выходное реле \emph{знает}, что соединение было отклонено, но выводит причину из другого значения. Соединения к назначениям неблокирующие: \texttt{connect()} сначала возвращает \texttt{EINPROGRESS} (errno 115), а реальный результат становится доступен позже через событие записи. В этот момент \file{connection.c} получает настоящую ошибку через \texttt{getsockopt(SO\_ERROR)}, записывает \enquote{Connection refused} и вызывает \file{connection\_edge\_end\_errno()}. Однако эта функция повторно вычисляет причину через \texttt{errno\_to\_stream\_end\_reason(tor\_socket\_errno(\ldots))}. В Unix \texttt{tor\_socket\_errno} определён как макрос \texttt{(errno)} (\file{src/lib/net/socket.h}), то есть обращается к глобальной переменной, где всё ещё остаётся \texttt{EINPROGRESS}. Значение 115 отсутствует в списке \file{src/core/or/reasons.c}, поэтому ветвь по умолчанию возвращает \texttt{MISC}.

Правильная информация существует, но не используется. Две ветви кода в этом участке тоже различаются: ветвь соединений между реле передаёт полученную ошибку дальше, а ветвь потоков её игнорирует. Согласно имеющимся отчёту и патчу в \file{docs/it/segnalazioni/}, в предыдущем испытании на частной сети с tor 0.4.9.13 передача настоящего кода ошибки изменила третий ответ на \texttt{CONNECTREFUSED}. Это небольшой дефект без последствий для анонимности: клиент всё равно знает, что поток завершился ошибкой, но теряет полезную информацию. \texttt{CONNECTREFUSED} показал бы, что выходное реле достигло назначения и повтор через другой exit бессмысленен. Анализ и патч находятся в \file{docs/it/segnalazioni/} для Tor Project. Этот пример также показывает, почему спецификации нужно проверять по реальному поведению.
\end{approfondimento}

Реальный локализованный скрипт — \file{labs/ru/лаборатория/exit\_policy.py}. В Linux с запущенной выделенной сетью Chutney его можно запустить из корня репозитория командой \texttt{PYTHONPATH=labs/ru/код python3 labs/ru/лаборатория/exit\_policy.py}. Он выполняет три проверки и пытается восстановить исходный \file{torrc} узла \texttt{006e} в блоке \texttt{finally}. Используйте только одноразовую лабораторную сеть: скрипт временно меняет настройки выходного реле и открывает локальные HTTP-серверы. Обычная цель \texttt{make -C labs/ru/лаборатория verify} проверяет только Chutney и не запускает скрипт. Его запуск и фактическое восстановление не проверялись в ходе аудита. Важная деталь: цепь должна содержать три перехода. Выходное реле отказывается пересылать трафик от одношаговой цепи и вернёт \texttt{END} с причиной \texttt{TORPROTOCOL} (\cref{sec:dos-limiti}).

% ---------------------------------------------------------------------------
\section{Семейства реле}\label{sec:relay-famiglie}

Операторы, управляющие несколькими реле, должны объявлять \textbf{семейство}. Клиенты никогда не помещают два реле из одного семейства в одну цепь (\cref{sec:percorso-vincoli}), поскольку общий оператор мог бы коррелировать трафик. Исторический механизм — \opt{MyFamily}: каждое реле перечисляет отпечатки всех остальных, и принадлежность к семейству учитывается только при взаимных объявлениях. Поэтому при десяти реле приходится поддерживать десять совпадающих списков, каждый из которых нужно менять при добавлении нового узла.

Механизм, введённый в tor 0.4.9, вместо этого использует \textbf{общий ключ семейства}. Оператор создаёт его командой \cmd{tor --keygen-family MyFamilyKey}, копирует файл \texttt{MyFamilyKey.secret\_family\_key} в каталог \texttt{KeyDirectory} каждого реле, сохраняя суффикс \texttt{.secret\_family\_key}, и добавляет в каждый \file{torrc} полученную строку \texttt{FamilyId}. Общий секретный ключ нужно защищать и хранить в безопасной резервной копии; при утечке его необходимо заменить на всех реле семейства. Каждое реле подписывает этим ключом сертификат принадлежности к семейству. Проверяющим больше не нужно сравнивать списки. На переходном этапе рекомендуется настраивать оба механизма.

\begin{approfondimento}[Почему объявление семейства полезно]
Ничто технически не заставляет оператора объявлять свои реле как семейство, и атакующий может намеренно этого не делать. Однако авторитетные серверы каталога и добровольцы, наблюдающие сеть, ищут общие характеристики: одни и те же сети, версии, синхронные времена запуска, контактные данные. Необъявленное семейство может привлечь внимание, и его реле могут быть исключены из консенсуса. Декларация — часть модели доверия, на которой держится сеть.
\end{approfondimento}

% ---------------------------------------------------------------------------
\textbf{Не включайте мосты в публичные списки \texttt{MyFamily}:} это может раскрыть их связь с опубликованными реле. Руководство Tor прямо рекомендует не использовать эту опцию на мостах. Кроме того, \texttt{tor --keygen-family MyFamilyKey} может перезаписать существующие файлы с таким именем: выберите новый базовый идентификатор и защитите созданный секретный файл.

\section{Запуск моста}\label{sec:relay-bridge}

Мост с obfs4 можно настроить, как в Лабораторной работе~10 (\cref{sec:lab-censura}):
\begin{lstlisting}[numbers=none]
BridgeRelay 1
ORPort 9001
ServerTransportPlugin obfs4 exec /usr/bin/obfs4proxy
ServerTransportListenAddr obfs4 0.0.0.0:443
ExtORPort auto
BridgeDistribution any
\end{lstlisting}
\opt{BridgeDistribution} сообщает системе распределения, сейчас rdsys, каким способом оператор хочет распространять мост. Значение \texttt{any} оставляет выбор системе, а \texttt{none} полностью исключает его из распространения, что удобно для передачи моста доверенным людям. Доступные каналы, включая веб-сайт, электронную почту и Telegram, перечислены на странице информации о мостах. Мост отсутствует в публичном консенсусе, но его дескриптор всё равно отправляется авторитетному серверу мостов, который проверяет доступность.

Подключаемый транспорт должен прослушивать порт, отличный от ORPort, и по возможности привычный порт вроде 443. Именно этот порт используют клиенты и видит цензор (\cref{sec:cens-strategie}).

% ---------------------------------------------------------------------------
\section*{Итоги}
\begin{itemize}
  \item Промежуточное реле — наименее требовательный способ помочь сети; выходное реле особенно ценно, но несёт больше внешних рисков; мост помогает пользователям в условиях цензуры.
  \item Существенно всего несколько параметров: идентичность, контактные данные, порты и лимиты пропускной способности. Основной ключ идентичности можно хранить вне сети.
  \item Политика выхода — упорядоченный список, где побеждает первое совпавшее правило; приватные назначения по умолчанию отклоняются.
  \item Отказ по политике приходит клиенту как сообщение \texttt{END} с причиной \texttt{EXITPOLICY}, адресом и TTL.
  \item Операторы нескольких реле должны объявлять их как семейство с помощью \opt{MyFamily} или более нового механизма общего ключа семейства.
\end{itemize}

% ---------------------------------------------------------------------------
\section*{Упражнения}

\begin{esercizio}[1]\label{es:22-ordine}
Для policy \texttt{accept *:80}, \texttt{reject 192.0.2.0/24:*}, \texttt{reject *:*} может ли клиент подключиться через этот exit к \texttt{192.0.2.7:80}? Перепишите policy так, чтобы это было невозможно.
\end{esercizio}

\begin{esercizio}[1]\label{es:22-riassunto}
Почему сводки только по портам в консенсусе недостаточно, чтобы определить, примет ли exit соединение? Что делает клиент после получения \texttt{END} с причиной \texttt{EXITPOLICY}?
\end{esercizio}

\begin{esercizio}[2]\label{es:22-famiglia}
Оператор управляет пятью relay через \opt{MyFamily}, затем добавляет шестой, но обновляет конфигурацию только на новом узле. Что видят клиенты? Что меняется при использовании family key?
\end{esercizio}

\begin{esercizio}[2]\label{es:22-refused}
Проверьте объяснение из технического разбора: включите \texttt{info}-журналы на exit и повторите три запроса. Что изменилось бы, если \file{connection\_edge\_end\_errno()} получала ошибку, возвращённую \texttt{getsockopt}? Какие ещё причины закрытия stream тогда могли бы сообщаться корректно?
\end{esercizio}

\begin{esercizio}[3]\label{es:22-accounting}
Для реле задано \texttt{AccountingMax 1 TBytes} в месяц и пропускная способность 10\,MB/s. Через сколько дней работы на полной скорости закончится лимит? Как реле, отключённое половину месяца, влияет на сеть по сравнению с узлом, работающим весь месяц на половине пропускной способности?
\end{esercizio}

% ---------------------------------------------------------------------------
\section*{Дополнительная литература}
\begin{itemize}
  \item Руководство Tor Project для операторов relay: \url{https://community.torproject.org/relay/}.
  \item \begin{sloppypar}Разделы руководства tor для \opt{ExitPolicy}, \opt{ReducedExitPolicy},
        \opt{FamilyId}, \opt{OfflineMasterKey} и \opt{BridgeDistribution}.\end{sloppypar}
  \item Спецификация причин закрытия stream: \spec{tor-spec}{Closing streams}.
\end{itemize}
